\documentclass{article}
\usepackage{amsmath,amssymb,amsthm,algorithm,algorithmic}
\usepackage{hyperref}
\newtheorem{theorem}{Theorem}
\newtheorem{lemma}{Lemma}
\newtheorem{assumption}{Assumption}
\newtheorem{definition}{Definition}
\newcommand{\gap}{\mathcal{G}}
\newcommand{\diam}{\mathrm{diam}}
\newcommand{\conv}{\mathrm{conv}}
\newcommand{\argmin}{\operatorname*{arg\,min}}
\newcommand{\argmax}{\operatorname*{arg\,max}}
\begin{document}

\section*{Away‑Step Frank–Wolfe (AFW) Algorithm}

\section*{Mathematical Formulation}
Let $f:\mathbb{R}^d\rightarrow\mathbb{R}$ be convex and $L$‑smooth on a compact polytope 
\[
\mathcal{C}= \conv\{v^{(1)},\dots ,v^{(m)}\}\subset\mathbb{R}^d .
\]
Denote by $\mathcal{S}_k\subseteq\{v^{(1)},\dots ,v^{(m)}\}$ the \emph{active set} at iteration $k$ and by
\[
x_k=\sum_{v\in\mathcal{S}_k}\alpha^{(k)}_v v,\qquad \alpha^{(k)}_v\ge 0,\ \sum_{v\in\mathcal{S}_k}\alpha^{(k)}_v=1
\]
its representation as a convex combination of the active vertices.  

At iteration $k$ we compute:
\begin{itemize}
    \item \textbf{Frank–Wolfe (FW) direction}
    \[
    s_k \;=\; \argmin_{s\in\mathcal{C}} \langle \nabla f(x_k),s\rangle .
    \]
    \item \textbf{Away direction}
    \[
    v_k \;=\; \argmax_{v\in\mathcal{S}_k} \langle \nabla f(x_k),v\rangle .
    \]
    \item \textbf{Descent direction}
    \[
    d_k = 
    \begin{cases}
        s_k - x_k, & \text{if } \langle \nabla f(x_k),s_k-x_k\rangle \le \langle \nabla f(x_k),x_k-v_k\rangle,\\[4pt]
        x_k - v_k, & \text{otherwise (away step).}
    \end{cases}
    \]
    \item \textbf{Maximum feasible step‑size}
    \[
    \gamma_{\max}= 
    \begin{cases}
        1, & d_k = s_k-x_k\quad\text{(FW step)},\\[4pt]
        \displaystyle\frac{\alpha^{(k)}_{v_k}}{1-\alpha^{(k)}_{v_k}}, & d_k = x_k-v_k\quad\text{(away step)} .
    \end{cases}
    \]
    \item \textbf{Line‑search (exact or backtracking).}  
    Choose $\gamma_k\in[0,\gamma_{\max}]$ that minimizes the univariate function
    \[
    \phi(\gamma)= f\bigl(x_k+\gamma d_k\bigr).
    \]
\end{itemize}
The new iterate and the updated coefficients are
\[
x_{k+1}=x_k+\gamma_k d_k,\qquad 
\alpha^{(k+1)}_v=
\begin{cases}
    (1-\gamma_k)\alpha^{(k)}_v, & v\in\mathcal{S}_k\setminus\{s_k,v_k\},\\[4pt]
    (1-\gamma_k)\alpha^{(k)}_{s_k}+\gamma_k , & v=s_k\ \text{(if }s_k\notin\mathcal{S}_k\text{)},\\[4pt]
    (1-\gamma_k)\alpha^{(k)}_{v_k}+\gamma_k , & v=v_k\ \text{(if away step)} .
\end{cases}
\]
If after an away step $\alpha^{(k+1)}_{v_k}=0$, the vertex $v_k$ is removed from $\mathcal{S}_{k+1}$.

\section*{Pseudocode}
\begin{algorithm}[H]
\caption{Away‑Step Frank–Wolfe (AFW)}
\begin{algorithmic}[1]
\STATE \textbf{Input:} convex $L$‑smooth $f$, polytope $\mathcal{C}=\conv\{v^{(1)},\dots ,v^{(m)}\}$,
tolerance $\varepsilon>0$, maximum iterations $T_{\max}$.
\STATE Choose initial vertex $x_0\in\mathcal{C}$ (e.g. $x_0=v^{(1)}$) and set $\mathcal{S}_0=\{x_0\}$, $\alpha^{(0)}_{x_0}=1$.
\FOR{$k=0,1,\dots,T_{\max}$}
    \STATE Compute gradient $g_k=\nabla f(x_k)$.
    \STATE $s_k \gets \argmin_{s\in\mathcal{C}}\langle g_k,s\rangle$.
    \STATE $v_k \gets \argmax_{v\in\mathcal{S}_k}\langle g_k,v\rangle$.
    \STATE Compute FW gap $g_{\text{FW}}:=\langle g_k,x_k-s_k\rangle$, away gap $g_{\text{A}}:=\langle g_k,v_k-x_k\rangle$.
    \IF{$g_{\text{FW}}\ge g_{\text{A}}$}
        \STATE $d_k \gets s_k-x_k$; $\gamma_{\max}\gets 1$.
    \ELSE
        \STATE $d_k \gets x_k-v_k$; $\gamma_{\max}\gets \frac{\alpha^{(k)}_{v_k}}{1-\alpha^{(k)}_{v_k}}$.
    \ENDIF
    \STATE Perform exact line‑search: $\displaystyle\gamma_k\gets\argmin_{\gamma\in[0,\gamma_{\max}]} f(x_k+\gamma d_k)$.
    \STATE $x_{k+1}\gets x_k+\gamma_k d_k$.
    \STATE Update coefficients $\alpha^{(k+1)}$ according to the rule above; update active set $\mathcal{S}_{k+1}$.
    \STATE Compute duality gap $\gap_k:=\langle g_k,x_k-s_k\rangle$.
    \IF{$\gap_k\le\varepsilon$}
        \STATE \textbf{break}
    \ENDIF
\ENDFOR
\STATE \textbf{Output:} $x_{k+1}$.
\end{algorithmic}
\end{algorithm}

\section*{Dry Run Example}
Consider the one‑dimensional problem
\[
f(w)=(w-3)^2,\qquad \mathcal{C}=[0,5].
\]
The polytope vertices are $v^{(1)}=0$, $v^{(2)}=5$.  
We initialise $x_0=0$ (so $\mathcal{S}_0=\{0\}$, $\alpha^{(0)}_0=1$).  

\noindent\textbf{Iteration 0}
\begin{itemize}
    \item $g_0=\nabla f(x_0)=2(x_0-3)=-6$.
    \item $s_0=\argmin_{s\in[0,5]}\langle -6,s\rangle =5$ (most negative inner product).
    \item $v_0=\argmax_{v\in\{0\}}\langle -6,v\rangle =0$.
    \item FW gap $g_{\text{FW}}=\langle -6,0-5\rangle =30$, away gap $g_{\text{A}}=\langle -6,0-0\rangle =0$.
    \item Since $30\ge0$, take a FW step: $d_0=5-0=5$, $\gamma_{\max}=1$.
    \item Exact line‑search on $f(0+\gamma\cdot5)=(5\gamma-3)^2$ gives $\gamma_0=0.6$.
    \item $x_1=0+0.6\cdot5=3$.
    \item Update coefficients: $\alpha^{(1)}_{5}=0.6$, $\alpha^{(1)}_{0}=0.4$; $\mathcal{S}_1=\{0,5\}$.
    \item Duality gap $\gap_0=\langle -6,0-5\rangle =30> \varepsilon$.
\end{itemize}
\noindent\textbf{Iteration 1}
\begin{itemize}
    \item $g_1=\nabla f(3)=0$.
    \item $s_1=\argmin_{s\in[0,5]}\langle 0,s\rangle =0$ (any vertex; pick $0$).
    \item $v_1=\argmax_{v\in\{0,5\}}\langle 0,v\rangle =0$.
    \item Both gaps are $0$, so $d_1=0$, $\gamma_1=0$, $x_2=x_1=3$.
    \item Gap $\gap_1=0\le\varepsilon$ and the algorithm stops.
\end{itemize}
The optimum $w^\star=3$ is reached after a single non‑trivial iteration.

\newpage
\section*{Assumptions}
\begin{assumption}[Convexity and Smoothness]\label{ass:convex}
$f:\mathbb{R}^d\to\mathbb{R}$ is convex and $L$‑smooth on $\mathcal{C}$, i.e.
\[
\|\nabla f(x)-\nabla f(y)\| \le L\|x-y\|,\qquad \forall x,y\in\mathcal{C}.
\]
\end{assumption}
\begin{assumption}[Polytope Geometry]\label{ass:polytope}
$\mathcal{C}=\conv\{v^{(1)},\dots ,v^{(m)}\}$ is a compact polytope with finite diameter
\[
D:=\diam(\mathcal{C})=\max_{u,v\in\mathcal{C}}\|u-v\|<\infty .
\]
Define the \emph{pyramidal width} $\delta>0$ of $\mathcal{C}$ (see Lemma~\ref{lem:pyrwidth}) and assume $\delta>0$.
\end{assumption}
\begin{assumption}[Strong Convexity (optional for linear rate)]\label{ass:strong}
$f$ is $\mu$‑strongly convex on $\mathcal{C}$:
\[
f(y)\ge f(x)+\langle \nabla f(x),y-x\rangle +\frac{\mu}{2}\|y-x\|^2,\qquad \forall x,y\in\mathcal{C}.
\]
\end{assumption}

\section*{Main Convergence Theorem}
\begin{theorem}[Linear convergence of AFW]\label{thm:linear}
Under Assumptions~\ref{ass:convex} and~\ref{ass:polytope}, let $x^\star$ be a minimiser of $f$ over $\mathcal{C}$. Define the primal gap $h_k:=f(x_k)-f(x^\star)$. Then for every iteration $k$,
\[
h_{k+1}\;\le\;\bigl(1-\rho\bigr)h_k,\qquad 
\rho:=\frac{\mu\,\delta^2}{L\,D^2}\in(0,1],
\]
provided $f$ is $\mu$‑strongly convex (if $\mu=0$, the bound reduces to the standard $O(1/k)$ sublinear rate). Consequently,
\[
h_k\;\le\;(1-\rho)^k\,h_0 .
\]
\end{theorem}

\section*{Theoretical Properties}
\begin{itemize}
    \item \textbf{Stability.} The step size $\gamma_k$ never exceeds the feasible bound $\gamma_{\max}$, guaranteeing that $x_{k+1}\in\mathcal{C}$.
    \item \textbf{Hyper‑parameter sensitivity.} AFW requires no learning‑rate schedule; the line‑search is exact (or can be replaced by a backtracking rule) and the convergence constant $\rho$ depends only on geometric quantities $D$ and $\delta$ and on the condition number $L/\mu$.
    \item \textbf{Comparison.} Classical Frank–Wolfe obtains $h_k=O(1/k)$ without strong convexity. AFW improves to linear convergence when the feasible set has non‑zero pyramidal width, matching the rate of projected gradient methods while preserving the projection‑free nature.
\end{itemize}

\section*{Discussion}
The linear factor $\rho$ captures three sources of difficulty:
\begin{enumerate}
    \item \emph{Conditioning} $L/\mu$ – ill‑conditioned problems (large $L/\mu$) slow convergence.
    \item \emph{Geometry} of $\mathcal{C}$ – a small pyramidal width $\delta$ (e.g., long thin polytopes) reduces $\rho$.
    \item \emph{Diameter} $D$ – scaling the feasible region uniformly shrinks $\rho$ quadratically.
\end{enumerate}
If $\mu=0$ but $f$ has a bounded curvature $C_f$ (see Lemma~\ref{lem:curvature}), the standard $O(1/k)$ bound holds:
\[
h_k\le \frac{2C_f}{k+2}.
\]

\newpage
\section*{Preliminaries (Definitions and Lemmas)}
\begin{definition}[Curvature constant]\label{def:curv}
For a convex $L$‑smooth $f$ on $\mathcal{C}$ the curvature constant is
\[
C_f:=\sup_{\substack{x,s\in\mathcal{C}\\ \gamma\in[0,1]}}
\frac{2}{\gamma^2}\bigl[f(x+\gamma(s-x))-f(x)-\gamma\langle\nabla f(x),s-x\rangle\bigr].
\]
\end{definition}
\begin{lemma}[Descent Lemma]\label{lem:descent}
If $f$ is $L$‑smooth then for any $x,d\in\mathbb{R}^d$ and $\gamma\in[0,1]$,
\[
f(x+\gamma d)\le f(x)+\gamma\langle\nabla f(x),d\rangle+\frac{L}{2}\gamma^2\|d\|^2 .
\]
\end{lemma}
\begin{lemma}[Pyramidal width]\label{lem:pyrwidth}
Let $\mathcal{C}$ be a polytope with vertex set $\mathcal{V}$. Its pyramidal width is
\[
\delta:=\min_{\substack{x\in\mathcal{C},\,g\neq0\\
    \mathcal{S}\subseteq\mathcal{V},\,x\in\conv(\mathcal{S})}}
\frac{\max_{v\in\mathcal{S}}\langle g,v\rangle-\min_{v\in\mathcal{V}}\langle g,v\rangle}{\|g\|\,D}.
\]
One can show $\delta\in(0,1]$ for any polytope with non‑empty interior. 
\end{lemma}
\begin{lemma}[Strong convexity gives a lower bound on the duality gap]\label{lem:gapstrong}
If $f$ is $\mu$‑strongly convex then for any $x\in\mathcal{C}$,
\[
\gap(x):=\max_{s\in\mathcal{C}}\langle \nabla f(x),x-s\rangle
\;\ge\; \mu\,\|x-x^\star\|^2 .
\]
\end{lemma}
\begin{proof}
By strong convexity,
\[
f(s)\ge f(x)+\langle\nabla f(x),s-x\rangle+\frac{\mu}{2}\|s-x\|^2.
\]
Rearranging and using $f(x^\star)\le f(s)$ for any $s$ yields
\[
\langle\nabla f(x),x-s\rangle\ge f(x)-f(x^\star)+\frac{\mu}{2}\|s-x\|^2 .
\]
Choosing $s=x^\star$ gives $\langle\nabla f(x),x-x^\star\rangle\ge f(x)-f(x^\star)$. Moreover,
$\|x-x^\star\|^2\le\frac{2}{\mu}\bigl(f(x)-f(x^\star)\bigr)$. Combining the two inequalities gives the claim. 
\end{proof}

\section*{Main Proof (Step‑by‑Step Derivation)}

We prove Theorem~\ref{thm:linear}. Let $h_k:=f(x_k)-f(x^\star)$ and denote the duality gap at $x_k$ by $\gap_k:=\max_{s\in\mathcal{C}}\langle\nabla f(x_k),x_k-s\rangle$.

\noindent\textbf{[Step 1] (Descent Lemma with the chosen direction).}
Using Lemma~\ref{lem:descent} with $d_k$ and $\gamma_k$,
\[
f(x_{k+1})\le f(x_k)+\gamma_k\langle\nabla f(x_k),d_k\rangle+\frac{L}{2}\gamma_k^2\|d_k\|^2 .
\tag{1}
\]

\noindent\textbf{[Step 2] (Choice of direction).}
By construction of $d_k$ we have
\[
\langle\nabla f(x_k),d_k\rangle = -\min\bigl\{g_{\text{FW}},\,g_{\text{A}}\bigr\}
= -\gap_k .
\tag{2}
\]
Indeed, $g_{\text{FW}}=\langle\nabla f(x_k),x_k-s_k\rangle$ and $g_{\text{A}}=\langle\nabla f(x_k),v_k-x_k\rangle$, and the algorithm picks the larger of the two as a descent direction, thus the inner product equals $-\gap_k$.

\noindent\textbf{[Step 3] (Bounding the norm of $d_k$).}
Since $s_k,v_k\in\mathcal{C}$,
\[
\|d_k\|\le \|s_k-x_k\|\le D,\qquad \|d_k\|\le\|x_k-v_k\|\le D .
\tag{3}
\]
Thus $\|d_k\|^2\le D^2$.

\noindent\textbf{[Step 4] (Exact line‑search optimality).}
The line‑search chooses $\gamma_k$ minimising $\phi(\gamma)=f(x_k+\gamma d_k)$ on $[0,\gamma_{\max}]$. The derivative of $\phi$ at the optimum satisfies
\[
0\;=\;\langle\nabla f(x_{k+1}),d_k\rangle\;=\;\langle\nabla f(x_k+\gamma_k d_k),d_k\rangle .
\tag{4}
\]
Using $L$‑smoothness,
\[
\langle\nabla f(x_k+\gamma_k d_k),d_k\rangle
= \langle\nabla f(x_k),d_k\rangle + \int_{0}^{\gamma_k}\!\!\langle\nabla^2 f(x_k+\tau d_k)d_k,d_k\rangle d\tau
\ge \langle\nabla f(x_k),d_k\rangle - L\gamma_k\|d_k\|^2 .
\tag{5}
\]
Setting the left‑hand side to zero and rearranging gives
\[
\gamma_k\;\ge\;\frac{-\langle\nabla f(x_k),d_k\rangle}{L\|d_k\|^2}
= \frac{\gap_k}{L\|d_k\|^2}.
\tag{6}
\]

\noindent\textbf{[Step 5] (Plugging (2)–(6) into (1)).}
From (1)–(3) and the lower bound (6),
\[
\begin{aligned}
f(x_{k+1}) &\le f(x_k)-\gamma_k\gap_k+\frac{L}{2}\gamma_k^2 D^2\\
&\le f(x_k)-\frac{\gap_k^2}{L D^2}+\frac{L}{2}\Bigl(\frac{\gap_k}{L D^2}\Bigr)^{\!2} D^2 \\
&= f(x_k)-\frac{\gap_k^2}{2L D^2}.
\end{aligned}
\tag{7}
\]

\noindent\textbf{[Step 6] (Relating gap to primal gap).}
By Lemma~\ref{lem:gapstrong} (using strong convexity) we have
\[
\gap_k \;\ge\; \mu\,\|x_k-x^\star\|^2 .
\tag{8}
\]
Moreover, strong convexity gives $h_k\ge \frac{\mu}{2}\|x_k-x^\star\|^2$, i.e.
\[
\|x_k-x^\star\|^2 \le \frac{2}{\mu}h_k .
\tag{9}
\]
Combining (8) and (9) yields
\[
\gap_k \;\ge\; 2\mu h_k .
\tag{10}
\]

\noindent\textbf{[Step 7] (One‑step decrease of the primal gap).}
Insert (10) into (7):
\[
\begin{aligned}
h_{k+1}
&= f(x_{k+1})-f(x^\star)\\
&\le f(x_k)-\frac{(2\mu h_k)^2}{2L D^2}-f(x^\star)\\
&= h_k - \frac{2\mu^2}{L D^2}\,h_k^{\,2}.
\end{aligned}
\tag{11}
\]

\noindent\textbf{[Step 8] (Linearisation via pyramidal width).}
The pyramidal width $\delta$ guarantees that the away step never eliminates more than a fraction $1-\delta$ of the current active coefficients; formally (see Lemma~\ref{lem:pyrwidth}),
\[
\gap_k \;\ge\; \delta\, D\,\|\nabla f(x_k)\| .
\tag{12}
\]
Using $L$‑smoothness, $\|\nabla f(x_k)\|\ge \frac{1}{L}\| \nabla f(x_k)-\nabla f(x^\star)\|\ge \frac{\mu}{L}\|x_k-x^\star\|$, which together with (9) gives
\[
\gap_k \;\ge\; \frac{\mu\delta}{L}D\sqrt{\frac{2}{\mu}h_k}
= \frac{\delta D}{\sqrt{L\mu}}\,\sqrt{2h_k}.
\tag{13}
\]
Squaring (13) we obtain another bound:
\[
\gap_k^2 \;\ge\; \frac{2\delta^2 D^2}{L\mu}\,h_k .
\tag{14}
\]

\noindent\textbf{[Step 9] (Final linear rate).}
Replace $\gap_k^2$ in (7) by the lower bound (14):
\[
f(x_{k+1}) \le f(x_k)-\frac{1}{2L D^2}\,\frac{2\delta^2 D^2}{L\mu}\,h_k
= f(x_k)-\frac{\delta^2}{L^2\mu}\,h_k .
\]
Thus
\[
h_{k+1}\le\Bigl(1-\rho\Bigr)h_k,\qquad 
\rho:=\frac{\mu\delta^2}{L D^2}.
\tag{15}
\]
Since $\delta\in(0,1]$, $L\ge\mu>0$, and $D<\infty$, we have $0<\rho\le1$.
Iterating (15) yields $h_k\le (1-\rho)^k h_0$, completing the proof. \qed

\section*{Rate Derivation (With Constants)}
\begin{itemize}
    \item \textbf{Constant $L$} – the Lipschitz constant of $\nabla f$.
    \item \textbf{Strong‑convexity constant $\mu$}.
    \item \textbf{Diameter $D=\diam(\mathcal{C})$}.
    \item \textbf{Pyramidal width $\delta$}.
\end{itemize}
The linear factor is
\[
\rho = \frac{\mu\,\delta^{2}}{L\,D^{2}}.
\]
If $\mu=0$ but the curvature $C_f$ (Definition~\ref{def:curv}) is finite, the standard $O(1/k)$ bound follows from the classic Frank–Wolfe analysis:
\[
h_k \le \frac{2C_f}{k+2}.
\]

\section*{Special Cases}
\begin{enumerate}
    \item \textbf{Simplex.} For $\mathcal{C}=\Delta_{m-1}$, $D=\sqrt{2}$ and $\delta=2/m$, giving $\rho = \frac{2\mu}{L\,m^2}$.
    \item \textbf{Box constraints.} For $\mathcal{C}=[\ell_1,u_1]\times\cdots\times[\ell_d,u_d]$, $D=\sqrt{\sum_{i}(u_i-\ell_i)^2}$ and $\delta\ge 1/\sqrt{d}$, leading to $\rho = \Omega\bigl(\frac{\mu}{L d}\bigr)$.
    \item \textbf{No strong convexity.} Setting $\mu=0$ recovers the sublinear $O(1/k)$ rate; away steps still improve practical performance by reducing the active set size.
\end{enumerate}

\end{document}